%% kapa: exam=1 round=30 year=2019 period=1 form=A subject=economics subject_ko=경제학원론 seq=7 number=47
%% source: Q-Net 감정평가사 기출문제 (cst00309) artlSeq=5205205 file=extracted/1008/5205205/2019년 제30회 감정평가사 1차시험 문제지/2019년 제30회 감정평가사 1차시험 문제지/2019년도 제30회 감정평가사 1차 1교시 A형.hwp
%% conversion: kordoc 4.24.51 -> tools/gen1.py (auto); review=proofread-sonnet; flags=none
\begin{question}{47}
기업 A와 B가 생산량 경쟁을 하는 시장수요곡선은 ${ P=\alpha-q_{ A } -q_{ B } }$로 주어졌다. 기업 A와 B는 동일한 재화를 생산하며, 평균비용은 ${ c }$로 일정하다. 기업 A의 목적은 이윤극대화이고, 기업 B의 목적은 손실을 보지 않는 범위 내에서 시장점유율을 극대화하는 것이다. 다음 설명 중 옳지 않은 것은? (단, ${ P }$는 시장가격, ${ q_{ A } }$는 기업 A의 생산량, ${ q_{ B } }$는 기업 B의 생산량이며, ${ c< \alpha }$이다.)
\begin{choices}
  \item 균형에서 시장가격은 $c$이다.
  \item 균형에서 기업 A의 이윤은 0보다 크다.
  \item 균형에서 기업 B의 이윤은 0이다.
  \item 균형에서 기업 B의 생산량이 기업 A보다 크다.
  \item 균형은 하나만 존재한다.
\end{choices}
\end{question}
\answer{2}
